% 视觉校样：本文件用章、课层级验证栏目、图像、表格和公式的页面关系；
% 正式书稿可以把同一套接口映射到其他目录层级。
% 正式稿替换文字时保留接口调用方式，并在每个侧边对象之后收束环绕状态。
\chapter{开端}
\chaptervisual{\mdiImage\hspace{0.5em}章图像模块（正式稿替换）}
\chapterlead{章首页在这里交代材料的范围、阅读顺序和页面节奏。正式稿把这段文字替换为本章导言，图像或档案可以占据同一网格的上半部。}
\begin{center}
  \small\color{bookgray}\mdiCompass\hspace{0.35em}对象定位\hspace{1em}\mdiDocument\hspace{0.35em}材料进入\hspace{1em}\mdiTimeline\hspace{0.35em}论证展开\hspace{1em}\mdiTarget\hspace{0.35em}结果回收
\end{center}
\chapterabstract{本段用于放置章的范围说明。它交代材料的时间、来源和阅读次序，让章首页的空间承担导览功能。正式稿在这里说明本章从哪个对象起步，经过哪些材料，在哪一个问题上收束。导览段只负责定位，具体论证从课页开始。}

\section{章首页的标题位置}
\lessonlead{课首页在标题之后交代对象、材料和本课要完成的论证动作，长度控制在三至五行；它不承担结论，也不堆叠关键词。}

\subsection{正文页的起笔}

这是正文占位段。正式写作时，第一段直接进入对象，说明材料从哪里来、问题在什么层面成立，以及读者将看到哪一个关系发生变化。段落之间用自然的论证连接，不用一组相同的卡片制造节奏。

第二段展示双面教材的行长和首行缩进。内侧页边留出装订空间，外侧页边放置页码和呼吸；页面底部保留比顶部更深的空白，正文因此不会贴近裁切线。
% 短引文必须从段首调用，正文在便利贴旁边环绕。
\sidequote{引文作者}{作品、版本、页码}{“这里粘贴经过版本核对的原文。”}
\noindent 引文旁边的正文沿便利贴的实际高度换行。这里保留足够长度的连续文字，用来检查小引文的真实环绕范围：正文从引文旁边开始排，经过数行后回到完整的版心，便利贴的高度不会制造一块额外的空白区域。正式稿中，小引文只在需要并置原文和论证时出现，出处始终留在便利贴内部。
\bookwrapclear

\begin{readingbox}
栏目样张。阅读与思考、问题引入和材料摘录各自有明确标题；底色只标出栏目类型，正文仍然是连续的阅读对象。
\end{readingbox}

\subsection{栏目的进入方式}

栏目应在论证需要停顿的位置出现。它可以给出一段原典、一道问题或一份材料，随后必须回到正文。栏目本身不代替论证，也不把装饰性的标签伪装成结论。

\begin{sourcebox}[原典与材料]
这里放置材料说明、版本差异或图像来源。直接引文使用上方的便利贴栏目，便利贴正文只保留原文和出处。
\end{sourcebox}

\relatedlink{这里放置与本节相邻的文本、图表或跨章提示。虚线只划出阅读路径，不把它做成一块新的信息卡片。}

\subsubsection{图片与正文的比例}

图像应当进入网格，而不是挤压正文。样张用占位块标出图像模块；正式稿替换为照片、档案扫描或原创图示，并在图下注明来源和说明。

\sideplaceholder[r]{0.36\textwidth}{34mm}{图 1　侧边图像与正文换行}
侧边图像用于材料高度有限、正文需要连续推进的段落。图像占据版心右侧，正文沿图像的实际高度换行；图像比例由原文件决定，排版接口不设裁剪高度。正式稿在这里替换为照片、档案扫描或原创图示，并在图下注明来源和说明。

图像之后，正文恢复为完整版心。编辑可以根据图像高度和段落长度选择左右位置；当图像过高，整块材料移到下一页，正文保持连续。
\bookwrapclear

\subsection{一页中的第二种节奏}

资料页可以让图像先出现，再让文字解释图像在本章中的作用。页面上的变化来自对象和论证的变化，颜色只承担导航。

\begin{questionbox}[问题引入]
问题框只提出一个可继续追踪的问题。问题之后紧跟材料或正文，避免把问句孤零零地放在页面上。
\end{questionbox}

% booktable 负责表格前后留白；表题与表格保持紧邻。
\tablecaption{表 1　页面元素与版式职责}
\begin{booktable}{@{}>{\bfseries\raggedright\arraybackslash}p{27mm}>{\raggedright\arraybackslash}X@{}}
\toprule
页面元素 & 版式职责\\
\midrule
章首页 & 建立章、课、对象的层级\\
正文页 & 承担连续论证与材料解释\\
栏目页 & 让读者停顿、核对或提出问题\\
图像模块 & 给出可指认的材料与来源\\
\bottomrule
\end{booktable}

这段文字用于检查表格之后的回到正文。正式稿在这里解释表格显示的关系，并把读者带回本节的主线。表格不承担结论，结论仍由连续文字完成。

\relatedlink{相关链接可以指向下一节的材料，也可以把读者送回本章开头的定义；它只提供位置，不制造新的理论层级。}

\subsection{材料回到正文}

侧边图和栏目结束后，正文恢复到完整版心。正式稿在这里继续处理材料的细节，说明图像、引文或表格怎样改变本课正在追踪的对象。

\begin{sourcebox}[材料说明]
这一小块用于登记图像的时间、地点、作者和版本。它只说明材料，解释由正文完成。
\end{sourcebox}

\bookwrapclear

\sidequote[l][0.24\linewidth]{引文作者}{作品、版本、页码}{“需要在续页保留的原文放在这里，正文沿引文之后恢复到版心。”}
\noindent 材料说明之后的段首保留一小段引导。这段文字用于检查材料说明之后的正文回收。这里用连续段落测试窄栏引文的环绕长度，正文在便利贴结束后恢复完整版心；正式稿在这里处理引文中的关键语句，并把出处、版本和页码留在便利贴内部。为使标题不进入便利贴的高度，环绕结束时由接口自动补足剩余行高。
\bookwrapclear

\subsection{引文之后的论证}

引文结束以后，正文回到同一版心，继续说明材料在本课中的位置。段落长度由论证决定，便利贴的高度随原文变化，页底保留给页码和装订安全区。

\begin{readingbox}[材料回看]
这里可以放置两至三句回看问题，指向引文中的具体词句。问题本身不替作者解释引文，答案留在后续正文。
\end{readingbox}

\begin{fullquote}{引文作者}{作品、版本、页码}
“这里放置需要在版心中完整阅读的较长引文。宽幅引文保留连续的正文尺度，折角只作为出处块的纸面标记。”
\end{fullquote}

\section{跨页阅读的样张}
\lessonlead{第二个课首页用于检查页眉、页码和标题在左右页的镜像关系。正式章节可以从这里开始一组新的材料，但不需要为了装饰额外插入空白页。}

\subsection{连续文字与页脚}

页眉只告诉读者此刻所在的章和课，页脚只承担页码。它们不抢正文的注意力，也不把作者说明写成正文。长段落在页底自然收束，孤立的一行会由分页规则自动避免。

\begin{noteBox}[制作提示]
这是给编辑和排版者看的工作提示。正式出版时可以删除，正文不应出现“写在开头”“本书将如何写作”一类元话语。
\end{noteBox}

\relatedlink{如果需要跨章回收概念，使用一句有对象的提示，指出读者应回到哪一段材料；不要用抽象口号替代交叉引用。}

\subsection{表格的留白}

表格只在比较关系确实比连续文字更清楚时使用。列宽服从正文版心，表头和表尾用细线，避免把每一个单元格都画成独立的框。

\subsection{公式的行长}

行内公式直接嵌入句子，例如 $f(x)=x^2+\lambda x$，其字面高度不改变中文正文的行距。需要单独阅读的推导使用带编号的行间公式：

% equation 产生按章递增的右侧编号；align 为每一行保留独立编号。
\begin{equation}
  \int_{0}^{T} f(t)\,\mathrm{d}t = F(T)-F(0)
  \label{eq:sample-integral}
\end{equation}

多行关系用 \texttt{align} 对齐，每一行可以单独交叉引用：

\begin{align}
  \Delta x &= x_{2}-x_{1}, \label{eq:sample-dx}\\
  \Delta y &= y_{2}-y_{1}. \label{eq:sample-dy}
\end{align}

正文可以写“由式 \eqref{eq:sample-integral} 得到……”，编号随章递增，行间距由母版统一控制。

本页结束后接入下一章。版心、页眉、页脚和栏目规则保持连续，章节文字由编辑稿替换。

\begin{readingbox}[课末回看]
课末回看把本课的对象、材料和论证动作重新并置。它可以是三句短问，也可以是一段回到原典的提示，长度由正文实际需要决定。
\end{readingbox}

\subsection{图文并置}

图像模块可以占据半个版心，说明文字紧随图像。正式稿在这里换入照片、档案或原创图示，并登记来源。

\fullplaceholder[0.46\textwidth]{第二图像模块（正式稿替换）}
\captionline{图 2　课内材料页的图像比例}
